\documentclass[11pt,a4paper,spanish]{book}
\usepackage{estilo_unir-1}
\usepackage{apacite}

\usepackage{amsmath}
\usepackage{amssymb}   

\let\rank\undefined % Forzamos que \rank quede libre
\usepackage{physics}
\usepackage{hyperref}
\numberwithin{equation}{chapter}
%\usepackage{caption}
\numberwithin{figure}{chapter}
%\usepackage{chngcntr}
%\counterwithin{equation}{chapter}
%\counterwithin{figure}{chapter}
%\counterwithin{table}{chapter}
%\renewcommand\theequation{\thechapter.\arabic{equation}}
%\renewcommand\thefigure{\thechapter.\arabic{figure}}
%\renewcommand\thetable{\thechapter.\arabic{table}}
%\renewcommand\theequation{\thechapter.\arabic{equation}}
%\counterwithin{figure}{chapter}
%\counterwithin{table}{chapter}
%\renewcommand\thefigure{\thechapter.\arabic{figure}}
%\renewcommand\thetable{\thechapter.\arabic{table}}
%\renewcommand\thefigure{\thechapter.\arabic{figure}}
%\makeatletter
%\renewcommand\p@figure{\thechapter..\arabic{figure}}
%\makeatother




%---------------------------
%título del trabajo y autor
%---------------------------
\title{Escribir el título del documento}
\titulacion{Máster en en Inteligencia Artificial}
\author{Gregory Harutyunyan Gevorgyan}
\date{\today}
\director{Pablo Abellan Galiana}
\nombreciudad{ Granada}

%---------------------------
%marges
%---------------------------
%\usepackage[margin=1.9cm]{geometry}
%---------------------------
%---------------------------
%---------------------------
%---------------------------
\begin{document}
\renewcommand{\listfigurename}{Índice de Ilustraciones}
\renewcommand{\listtablename}{Índice de Tablas}
\renewcommand{\contentsname}{Índice de Contenidos}
\renewcommand{\figurename}{Figura}
\renewcommand{\tablename}{Tabla} 

\maketitle
 
\frontmatter
\tableofcontents
\listoffigures
\listoftables

\chapter{Resumen}
{\bf Nota:} En este apartado se introducirá un breve resumen en español del trabajo realizado (extensión máxima: 150 palabras). Este resumen debe incluir el objetivo o propósito de la investigación, la metodología, los resultados y las conclusiones.


{\bf Palabras Clave:} Se deben incluir de 3 a 5 palabras claves en español

\chapter{Abstract}
{\bf Nota:} En este apartado se introducirá un breve resumen en español del trabajo realizado (extensión máxima: 150 palabras). Este resumen debe incluir el objetivo o propósito de la investigación, la metodología, los resultados y las conclusiones.


{\bf Palabras Clave:} Se deben incluir de 3 a 5 palabras claves en inglés




\mainmatter
\chapter{Introducción}\label{ch:introduccion}

\section{Motivación y Justificación}
% Aquí definiremos el problema de la decoherencia y el ruido en la era NISQ (Noisy Intermediate-Scale Quantum). 
% Justificaremos por qué la mitigación por IA es superior en coste computacional a la corrección de errores cuánticos (QEC) clásica.

\section{Planteamiento del Problema}
% Definición matemática del mapeo de un estado ruidoso a uno ideal.

\section{Estructura del Documento}
% Breve resumen de los capítulos posteriores.
\chapter{Contexto y Estado del Arte}\label{ch:estado_arte}


%\section{Fundamentos Matemáticos de Computación Cuántica}
%\subsection{Espacios de Hilbert y Cúbits}
% Definición del cúbit como vector unitario en \mathbb{C}^2 y sistemas multiqubit mediante productos tensoriales.
%\subsection{Puertas Cuánticas como Matrices Unitarias}
%\subsection{Canales Cuánticos y Modelos de Ruido}
% Formalismo de operadores de Kraus para modelar canales de despolarización, relajación térmica y error de lectura.

%\section{Mitigación de Ruido en Dispositivos NISQ}
%\subsection{Técnicas Clásicas de Mitigación (ZNE, PEC)}
%\subsection{Aproximación mediante Machine Learning}
% Revisión de papers actuales donde se usan redes neuronales para aproximar observables u operadores densidad.


\section{Fundamentos Matemáticos de Computación Cuántica}

\subsection{Espacios de Hilbert y Cúbits}

El estado de un sistema cuántico elemental, o \textit{cúbit}, se representa como un vector unitario $\ket{\psi}$ en el espacio de Hilbert complejo $\mathcal{H} = \mathbb{C}^2$. Utilizando la notación de Dirac, la base computacional ortonormal $\{\ket{0}, \ket{1}\}$ se define mediante los vectores de la base canónica:
\begin{equation}
\ket{0} = \begin{pmatrix} 1 \\ 0 \end{pmatrix}, \quad \ket{1} = \begin{pmatrix} 0 \\ 1 \end{pmatrix}
\end{equation}
Cualquier estado puro se expresa como una combinación lineal $\ket{\psi} = \alpha\ket{0} + \beta\ket{1}$, donde $\alpha, \beta \in \mathbb{C}$ satisfacen la condición de normalización $|\alpha|^2 + |\beta|^2 = 1$.

Para un sistema compuesto de $n$ cúbits, el espacio de estados resultante es el producto tensorial de los espacios individuales, $\mathcal{H}_{total} = \bigotimes_{i=1}^{n} \mathbb{C}^2 \cong \mathbb{C}^{2^n}$. Un estado general en este espacio de dimensión $2^n$ se escribe como:
\begin{equation}
\ket{\Psi} = \sum_{i=0}^{2^n-1} c_i \ket{i}, \quad \sum_{i=0}^{2^n-1} |c_i|^2 = 1
\end{equation}
donde $\ket{i}$ representa el producto tensorial de los estados de la base de cada cúbit.

\subsection{Puertas Cuánticas como Matrices Unitarias}

La evolución temporal de un sistema cuántico cerrado es reversible y preserva la norma, lo que implica que cualquier operación lógica o \textit{puerta cuántica} se describe matemáticamente como un operador unitario $U: \mathcal{H} \to \mathcal{H}$. Un operador $U$ es unitario si satisface $U^{\dagger}U = UU^{\dagger} = I$, donde $U^{\dagger}$ es la adjunta de la matriz (traspuesta conjugada).

En el contexto de ingeniería, esto implica que las puertas cuánticas son isometrías en la norma $L_2$. Para un sistema de $n$ cúbits, $U \in \mathbb{C}^{2^n \times 2^n}$. La fidelidad de una operación real frente a la ideal es una métrica crítica para el entrenamiento de modelos de IA de mitigación.

\subsection{Canales Cuánticos y Modelos de Ruido}

En dispositivos reales (NISQ), los sistemas no están aislados. La interacción con el entorno transforma estados puros en mezclas estadísticas, descritas por el operador densidad $\rho = \sum_i p_i \ket{\psi_i}\bra{\psi_i}$, donde $\rho$ es semidefinida positiva y $\text{Tr}(\rho) = 1$.

La evolución de un sistema abierto se modela mediante un \textit{canal cuántico} $\mathcal{E}$, un mapa completamente positivo y que preserva la traza (CPTP). Utilizando la representación de operadores de Kraus:
\begin{equation}
\mathcal{E}(\rho) = \sum_{k} E_k \rho E_k^{\dagger}, \quad \sum_{k} E_k^{\dagger} E_k = I
\end{equation}
Los modelos de ruido fundamentales en este TFM incluyen:
\begin{enumerate}
    \item \textbf{Canal de Despolarización:} Representa la pérdida de información del estado hacia la identidad: $\mathcal{E}(\rho) = (1-p)\rho + \frac{p}{2}I$.
    \item \textbf{Relajación Térmica ($T_1$ y $T_2$):} Describe la transición hacia el estado fundamental ($T_1$) y la pérdida de coherencia de fase ($T_2$).
    \item \textbf{Error de Lectura (ROME):} Modelado como una matriz de confusión clásica $P(m|i)$ que distorsiona la distribución de probabilidad tras la medida.
\end{enumerate}

\section{Mitigación de Ruido en Dispositivos NISQ}

\subsection{Técnicas Clásicas de Mitigación (ZNE, PEC)}

La mitigación de ruido, a diferencia de la corrección de errores (QECC), no utiliza cúbits redundantes, sino post-procesamiento estadístico para estimar el valor esperado ideal $\langle O \rangle_{ideal}$.

\begin{itemize}
    \item \textbf{Zero Noise Extrapolation (ZNE):} Consiste en escalar artificialmente el ruido del hardware por un factor $\lambda > 1$ (mediante el estiramiento de pulsos o inserción de puertas identidad) y ajustar una curva funcional para extrapolar el resultado al límite $\lambda \to 0$.
    \item \textbf{Probabilistic Error Cancellation (PEC):} Descompone el inverso del canal de ruido $\mathcal{E}^{-1}$ como una combinación lineal de operaciones implementables. Se muestrean circuitos de forma estocástica, lo que permite obtener un estimador insesgado de $\langle O \rangle$, a costa de un incremento exponencial en la varianza (muestreo).
\end{itemize}

\subsection{Aproximación mediante Machine Learning}
% Revisión de papers actuales donde se usan redes neuronales para aproximar observables u operadores densidad.
\chapter{Identificación de Requisitos}\label{ch:requisitos}

\section{Requisitos Funcionales del Sistema de IA}
% Definición de inputs (circuitos/observables ruidosos) y outputs (valores mitigados).

\section{Requisitos No Funcionales y Restricciones (NISQ)}
% Latencia de inferencia permitida, restricciones de memoria y dimensionalidad del espacio de Hilbert (que crece 2^n).

\section{Ecosistema Tecnológico y MLOps}
% Justificación analítica de las herramientas:
% - Qiskit Aer para simulación de modelos de ruido.
% - PyTorch para arquitecturas de Deep Learning.
% - MLflow / Weights & Biases para tracking de hiperparámetros y reproducibilidad.
\chapter{Objetivos}\label{ch:objetivos}

\section{Objetivo General}
% Desarrollar, entrenar y desplegar un modelo de Machine Learning capaz de mitigar errores de lectura y depolarización en circuitos parametrizados.

\section{Objetivos Específicos}
\subsection{Ingeniería de Datos Cuánticos}
\subsection{Diseño y Entrenamiento del Modelo}
\subsection{Implementación del Pipeline MLOps}
\subsection{Evaluación Out-Of-Distribution (OOD)}
\chapter{Desarrollo del Trabajo}\label{ch:desarrollo}

\section{Generación del Dataset Cuántico (Qiskit)}
\subsection{Definición de Circuitos Ansatz}
\subsection{Simulación de Operadores de Ruido}

\section{Arquitectura del Modelo de Aprendizaje (PyTorch)}
\subsection{Diseño de la Red}
% Trade-off entre capacidad del modelo (complejidad) y riesgo de sobreajuste.
\subsection{Función de Pérdida (Loss Function)}
% Justificación de MSE vs Fidelidad u otras distancias matriciales.

\section{Ciclo de Vida MLOps y Trazabilidad}
\subsection{Control de Semillas y Reproducibilidad}
\subsection{Tracking de Experimentos (MLflow/W\&B)}
% Análisis de las ejecuciones, Learning Rates y convergencia.

\section{Resultados y Evaluación Analítica}
\subsection{Rendimiento In-Distribution}
\subsection{Validación Out-Of-Distribution}
% Crucial: Demostrar que el modelo aprende la estructura del ruido y no memoriza los datos de entrenamiento.
\chapter{Conclusiones y Trabajo Futuro}

Aquí va la conclusión


\nocite{*}
%\begin{thebibliography}{a}
%\bibitem{etiqueta} \textsc{Autores},
%\textit{nombre referencia.}
%Información addicional
%\end{thebibliography}
\bibliographystyle{apacite}
\bibliography{bibliografia}

\appendix
\chapter{Apendices}

\end{document}





















